\paragraph{C.3aw (Proved/refuted): low restriction kernels in positive characteristic.}
For every prime $p$ and integers
\[
                        r\ge2,\qquad r+1\le N\le2r-1,
\]
the construction below specifies an excellent normal strictly henselian
local ring of dimension $N$, with algebraically closed coefficient and
residue field of characteristic $p$, and a principal divisor such that
both punctures $X,Y$ are regular and
\[
 \ker\!\left(
  \mathrm{CH}^r(X)_{\mathrm{alg}}\longrightarrow
  \mathrm{CH}^r(Y)_{\mathrm{alg}}\right)\otimes\mathbb Q\ne0.
\]
In fact the kernel contains a specified non-torsion principal Gysin
image of an actual ordinary class on $Y$. The same class also proves
that $\mathrm{CH}^r(X)_{\mathrm{alg}}\to
\mathrm{CH}^r(X\setminus Y)_{\mathrm{alg}}$ is not injective.
The coefficient field is explicitly transcendental over
$\overline{\mathbb F}_p$; no rational nonvanishing over that constant
field itself is asserted.

\begin{proof}
\emph{Step 1: the fixed non-torsion zero-cycle and its field.}
Put $m=N-r-1$, so $0\le m\le r-2$. Take the precisely specified
finite-field Calabi--Yau $r$-fold $T_0$ of C.3at, including its
Fermat/triangular parameter rule in every prime. Enlarge its finite
field of definition to contain the specified constant point $a$ of
C.3at, and call the resulting field $\mathbb F_q$. Define
\[
 k_0=\overline{\mathbb F}_p,
 \quad K=k_0(T_0),
 \quad k=\overline K,
 \quad T=T_0\times k,
 \quad H_T=c_1(\mathcal O_T(1)).
\]
Let $P$ be its geometric generic point, and set $\alpha=[P]-[a]$.
The slope-zero and two-support proof in C.3at gives
\begin{equation}\label{eq:C3aw-nonsupport}
 \alpha\ne0\quad\text{in}\quad
 \mathrm{CH}^r(T)_{\mathbb Q}/
                    H_T\mathrm{CH}^{r-1}(T)_{\mathbb Q}.
\end{equation}
The point $a$ and the generic point $P$ are already rational over
$\mathbb F_q(T_0)$, so the choices made next require no larger
transcendental coefficient field.

\emph{Step 2: a specified smooth connected base through both points.}
Choose an integer $\ell>0$ and $m+1$ sections
\[
 s_0,\ldots,s_m\in
 H^0\bigl(T_{K},\mathcal I_{\{P,a\}}\mathcal O_{T_K}(\ell)\bigr)
\]
whose common zero scheme $B$ is smooth of codimension $m+1$.
We prove existence and give a terminating choice rule.
For sufficiently large $\ell$, the system of sections vanishing at
the two points generates off them and defines an immersion there,
by \citet[Lemma 33.47.1]{D_StacksBertini0FD4}. The map to the two
first-jet spaces is surjective: apply Serre vanishing to the exact
sequence from the square of their ideal to their ideal
\citep[Lemma 30.17.1, Tag 0B5U]{D_StacksSerre0B5U}. Thus at
$P$ and $a$ one can prescribe $m+1$ independent differentials.
By \citet[Lemma 33.47.3 and its proof]{D_StacksBertini0FD4},
successive general members are smooth away from the two points;
the independent jets make them smooth at the two points too.
This gives a nonempty open set of good tuples. The final dimension is
\[
                            \dim B=r-m-1\ge1.
\]
The members may all be chosen in the same sufficiently large degree;
restriction of an immersion to each smooth intermediate section
preserves the Bertini hypotheses. Because $B$ is a positive-dimensional
complete intersection in projective space, its Koszul sequence and
projective-space cohomology give $H^0(B,\mathcal O_B)=k$ after base
change to $k$. Thus the smooth $B_k$ is connected and integral.

To make the choice explicit, form for each $\ell$ the lexicographic
monomial basis of $H^0(T_0,\mathcal O(\ell))$ modulo its homogeneous
equation and solve the two linear evaluation equations at $P,a$ by
reduced row echelon form over $\mathbb F_q(T_0)$. To fix the
evaluations, trivialize the invertible sheaf using the first nonzero
homogeneous coordinate of each point. This gives a fixed
basis of the displayed space. Enumerate finite extensions
$\mathbb F_{q^b}$ and the coefficient tuples in this basis, and
enumerate $(\ell,b)$ by increasing sum, breaking ties lexicographically.
Choose the first tuple for which the projective Jacobian rank test
gives a smooth complete intersection of codimension $m+1$.
All tests are finite polynomial operations over the stated function
field. This search terminates: the good tuples form a nonempty open
over $K$, and the constant tuples $k_0^t$ are Zariski dense in
$\mathbb A^t_K$. Indeed a polynomial over any extension of an infinite
field which vanishes on all tuples from that field is zero, by
induction on the number of variables. Every constant tuple lies in
some finite extension of $\mathbb F_q$. The chosen sections therefore
have coefficients in $K$, and all the following objects are defined
over the same $k=\overline K$.

In the rest of the proof write $B$ for its base change to $k$.
The class $[P]-[a]\in\mathrm{CH}_0(B)$ is ordinarily algebraically
trivial: use $B$ itself if it is a curve, and otherwise a smooth
connected sufficiently ample curve in $B$ through the two points.
This supplies an actual fixed-product proper-curve family, exactly as
in C.3at. Its pushforward to $T$ is the fixed class $\alpha$.

\emph{Step 3: the smooth incidence divisor and its ordinary Gysin image.}
Set
\[
 V=\mathbb P^m\times T,
 \quad L=\mathcal O_{\mathbb P^m}(1)\boxtimes\mathcal O_T(\ell),
 \quad h=c_1(\mathcal O_{\mathbb P^m}(1)),
\]
and define the divisor
\[
 j:Q=V\!\left(\sum_{i=0}^m x_i s_i\right)\hookrightarrow V.
\]
This is a section of $L$. Off $B$ its projection to $T$ is a
$\mathbb P^{m-1}$-bundle when $m>0$, so is smooth there. When
$m=0$, it has no points off $B$. Along $B$, the differentials
$ds_0,\ldots,ds_m$ are independent; hence
$\sum x_i ds_i\ne0$ at every projective point $[x_0:\cdots:x_m]$.
This proves smoothness of $Q$ in all cases. It contains
\[
                         D=\mathbb P^m\times B.
\]
Write $u:D\hookrightarrow Q$ for this closed immersion and
$v:D\to B$ for the projection. Define the integral ordinary class
\[
              \beta_Q=u_*v^*([P]-[a])
                       \in\mathrm{CH}^{r-1}(Q)_{\mathrm{alg}}.
\]
The codimension is correct: $\dim Q=m+r-1$, whereas each fibre
$\mathbb P^m\times\{P\}$ has dimension $m$. Proper pushforward
and flat pullback of the actual family on $B$ prove ordinary algebraic
triviality, without a formal-family argument. Its pushforward is
\[
                j_*\beta_Q=\operatorname{pr}_T^*\alpha
                                      \in\mathrm{CH}^r(V).
\]

\emph{Step 4: non-torsion in the cone quotient.}
The projective-bundle formula gives
$\mathrm{CH}^*(V)=\mathrm{CH}^*(T)[h]/(h^{m+1})$.
Since $c_1(L)=h+\ell H_T$, evaluation at $h=-\ell H_T$ gives
the exact quotient
\[
 \frac{\mathrm{CH}^r(V)_{\mathbb Q}}
      {c_1(L)\mathrm{CH}^{r-1}(V)_{\mathbb Q}}
 \simeq
 \frac{\mathrm{CH}^r(T)_{\mathbb Q}}
      {(\ell H_T)^{m+1}
                 \mathrm{CH}^{r-m-1}(T)_{\mathbb Q}}.
\]
To see the kernel directly, eliminate powers of $h$ successively
using $h=-\ell H_T$; the only remaining relation is
$(-\ell H_T)^{m+1}=0$ in degree $m+1$. The denominator on the
right lies inside $H_T\mathrm{CH}^{r-1}(T)_{\mathbb Q}$.
Thus \eqref{eq:C3aw-nonsupport} proves that
$\operatorname{pr}_T^*\alpha$ has infinite order in the integral
cone quotient. This includes $m=0$.

\emph{Step 5: strict local realization and both restriction kernels.}
Take the full section ring
\[
 A=\bigoplus_{n\ge0}H^0(V,L^n),
 \quad R=A_{A_+}^{\mathrm{sh}},
 \quad f=\sum_{i=0}^m x_i s_i\in A_1\subset R,
\]
at the geometric vertex with residue $k$.
The section ring is normal and generated in degree one. For generation,
the projectively normal hypersurface ring of $T$ has surjective
multiplication of sections in every degree; so does its $\ell$th
Veronese. Tensoring with the analogous multiplication on
$\mathbb P^m$ proves the assertion for the Segre section ring.
Normality also follows from the standard valuation description of the
full section ring of an invertible sheaf on a normal integral variety:
it is the intersection, inside $k(V)[t]$, of the inequalities given
by the prime divisors and the nonnegative grading. Each corresponding
valuation ring is integrally closed. In particular this is the full
normal section ring required by C.3e.

The punctured cone is the punctured line bundle $L^{-1}$ over $V$,
hence is smooth. Its local ring at the vertex is excellent and normal;
strict henselization retains normality and has residue $k$.
The ring $R$ has dimension $\dim V+1=m+r+1=N$ and its puncture
$X$ is regular. The element $f$ is a nonzero divisor. Away from
the vertex its zero scheme is the punctured line bundle over the
smooth divisor $Q$, so $Y=X\cap V(f)$ is regular too. No normality
or full-section-ring claim about the quotient at its closed point is
needed.

Let $i:Y\hookrightarrow X$ be the principal Cartier immersion.
Pull the actual class $\beta_Q$ first to the punctured bundle over
$Q$ and then to $Y$, and call the resulting class
$\beta\in\mathrm{CH}^{r-1}(Y)_{\mathrm{alg}}$.
Flat base change and proper pushforward give
\[
                   i_*\beta=\gamma,
\]
where $\gamma$ is the pullback of
$\operatorname{pr}_T^*\alpha$. The integral strict-cone injection
C.3e and Step~4 show that $\gamma$ has infinite order. These are
actual families on the fixed punctures, since both the flat pullbacks
and the proper pushforwards are performed on the product with the
same smooth parameter curve before taking its two fibres.

Finally $\mathcal O_X(Y)\simeq\mathcal O_X$ because $f$ globally
generates the Cartier ideal. Thus $N_{Y/X}\simeq\mathcal O_Y$ and
the principal self-intersection formula gives, integrally,
\[
             i^*\gamma=i^*i_*\beta
                         =c_1(N_{Y/X})\cap\beta=0.
\]
This is the all-characteristic identity proved in C.3al, with its
exact primary formulas in
\citet[Tags 0B3P, 02T8 and 0B6Z]{cyStacksPrincipalGysin}.
The class also restricts to zero on $X\setminus Y$ by localization,
since it is a closed pushforward from $Y$. Both asserted failures
therefore follow from the same non-torsion ordinary class.
\end{proof}

The proof leaves the positive-characteristic injectivity range
$N\ge2r$ undecided. For $N\le r$ the source has no codimension-$r$
points, so injectivity is automatic. The generic-point non-torsion
input uses the slope proof of C.3at and is not inferred from differential
forms or from a finite characteristic computation.
